% GENERATED FILE. Edit canonical lessons, then run npm run generate:books.
% canonical-source-hash: fnv1a64:f56211e2d3d1b038
% canonical-lessons: MR-C485-madhumeh, MR-C485-raktadab, MR-C485-dama, MR-C485-kavil, MR-C485-kanjinya

\chapter{Long Illnesses}
\label{ch:mr-long-illnesses}
\hlchaptermodality{485}

\begin{chapteropening}
\textbf{By the end of this chapter:} \emph{I can name diabetes, blood pressure, asthma, jaundice and chickenpox.}

Complete the last lesson of chapter 485: I can name diabetes, blood pressure, asthma, jaundice and chickenpox.
\end{chapteropening}

\section[\texorpdfstring{\mr{मधुमेह} (madhumeh)}{मधुमेह (madhumeh)}]{\mr{मधुमेह} (madhumeh) — diabetes}

\label{lesson:MR-C485-madhumeh}



\begin{quote}
Before the new one: say the Marathi for a virus, then the Marathi for an epidemic.
\end{quote}

\subsection*{You'll want to know: \mr{मधुमेह}}
\textbf{\mr{मधुमेह}} — \emph{madhumeh} — \textquotedblleft{}diabetes\textquotedblright{}.

\begin{quote}
\mr{मधुमेहामुळे} \mr{आजोबा} \mr{साखर} \mr{खात} \mr{नाहीत}.
\end{quote}

\begin{quote}
\emph{madhumehāmuḷe ājobā sākhar khāt nāhīt.}
\end{quote}

\begin{quote}
Grandfather does not eat sugar because of his diabetes.
\end{quote}

Masculine.

\mr{मधु}, honey, + \mr{मेह}, urine: an old name for the sugar sickness.

\begin{cousinweb}
\textbf{\mr{मधुमेह}} (\emph{madhumeh}) is a tatsama, taken from Sanskrit \emph{madhumeha}, \textquotedblleft{}honey-like urine, diabetes\textquotedblright{} (Monier-Williams).
\end{cousinweb}

\begin{grammarlens}[title={What you've built}]
One more word for describing what you lived and felt.
\end{grammarlens}

\subsection*{Your turn}
\begin{itemize}
\raggedright
  \item \emph{Say it:} \emph{madhumeh}
  \item \emph{Say it:} \emph{madhumeh}, once more
  \item \emph{Say it:} \emph{madhumehāmuḷe ājobā sākhar khāt nāhīt.} — Grandfather does not eat sugar because of his diabetes.
  \item \emph{Recall:} say the Marathi for a virus, then the Marathi for an epidemic, then say \emph{madhumeh} again
\end{itemize}

\begin{tcolorbox}[breakable,colback=teal!4,colframe=teal!35!black,title={Before you move on}]
What does \mr{मधुमेह} mean? (\textquotedblleft{}Diabetes\textquotedblright{}.) Say it once more.
\end{tcolorbox}
\section[\texorpdfstring{\mr{रक्तदाब} (raktadāb)}{रक्तदाब (raktadāb)}]{\mr{रक्तदाब} (raktadāb) — blood pressure}

\label{lesson:MR-C485-raktadab}



\begin{quote}
Before the new one: say the Marathi for an epidemic, then the Marathi for diabetes.
\end{quote}

\subsection*{You'll want to know: \mr{रक्तदाब}}
\textbf{\mr{रक्तदाब}} — \emph{raktadāb} — \textquotedblleft{}blood pressure\textquotedblright{}.

\begin{quote}
\mr{दवाखान्यात} \mr{त्यांनी} \mr{माझा} \mr{रक्तदाब} \mr{तपासला}.
\end{quote}

\begin{quote}
\emph{davākhānyāt tyānnī mājhā raktadāb tapāslā.}
\end{quote}

\begin{quote}
At the clinic they checked my blood pressure.
\end{quote}

Masculine.

\mr{रक्त} + \mr{दाब}, pressure, from \mr{दाबणे}.

\begin{cousinweb}
\textbf{\mr{रक्तदाब}} (\emph{raktadāb}) is a modern Marathi compound of \mr{रक्त}, \textquotedblleft{}blood\textquotedblright{}, a tatsama of Sanskrit \emph{rakta} (CDIAL 10539), and \mr{दाब}, \textquotedblleft{}pressure\textquotedblright{}, the noun of \mr{दाबणे}.
\end{cousinweb}

\begin{grammarlens}[title={What you've built}]
One more word for describing what you lived and felt.
\end{grammarlens}

\subsection*{Your turn}
\begin{itemize}
\raggedright
  \item \emph{Say it:} \emph{raktadāb}
  \item \emph{Say it:} \emph{raktadāb}, once more
  \item \emph{Say it:} \emph{davākhānyāt tyānnī mājhā raktadāb tapāslā.} — At the clinic they checked my blood pressure.
  \item \emph{Recall:} say the Marathi for an epidemic, then the Marathi for diabetes, then say \emph{raktadāb} again
\end{itemize}

\begin{tcolorbox}[breakable,colback=teal!4,colframe=teal!35!black,title={Before you move on}]
What does \mr{रक्तदाब} mean? (\textquotedblleft{}Blood pressure\textquotedblright{}.) Say it once more.
\end{tcolorbox}
\section[\texorpdfstring{\mr{दमा} (damā)}{दमा (damā)}]{\mr{दमा} (damā) — asthma}

\label{lesson:MR-C485-dama}



\begin{quote}
Before the new one: say the Marathi for diabetes, then the Marathi for blood pressure.
\end{quote}

\subsection*{You'll want to know: \mr{दमा}}
\textbf{\mr{दमा}} — \emph{damā} — \textquotedblleft{}asthma\textquotedblright{}.

\begin{quote}
\mr{धुळीमुळे} \mr{त्याचा} \mr{दमा} \mr{वाढतो}.
\end{quote}

\begin{quote}
\emph{dhuḷīmuḷe tyācā damā vāḍhto.}
\end{quote}

\begin{quote}
Dust makes his asthma worse.
\end{quote}

Masculine.

Masculine; \mr{दम} \mr{लागणे}, to get breathless.

\begin{cousinweb}
\textbf{\mr{दमा}} (\emph{damā}) is a loan, usually traced to Persian \emph{dam}, \textquotedblleft{}breath\textquotedblright{}.
\end{cousinweb}

\begin{grammarlens}[title={What you've built}]
One more word for describing what you lived and felt.
\end{grammarlens}

\subsection*{Your turn}
\begin{itemize}
\raggedright
  \item \emph{Say it:} \emph{damā}
  \item \emph{Say it:} \emph{damā}, once more
  \item \emph{Say it:} \emph{dhuḷīmuḷe tyācā damā vāḍhto.} — Dust makes his asthma worse.
  \item \emph{Recall:} say the Marathi for diabetes, then the Marathi for blood pressure, then say \emph{damā} again
\end{itemize}

\begin{tcolorbox}[breakable,colback=teal!4,colframe=teal!35!black,title={Before you move on}]
What does \mr{दमा} mean? (\textquotedblleft{}Asthma\textquotedblright{}.) Say it once more.
\end{tcolorbox}
\section[\texorpdfstring{\mr{कावीळ} (kāvīḷ)}{कावीळ (kāvīḷ)}]{\mr{कावीळ} (kāvīḷ) — jaundice}

\label{lesson:MR-C485-kavil}



\begin{quote}
Before the new one: say the Marathi for blood pressure, then the Marathi for asthma.
\end{quote}

\subsection*{You'll want to know: \mr{कावीळ}}
\textbf{\mr{कावीळ}} — \emph{kāvīḷ} — \textquotedblleft{}jaundice\textquotedblright{}.

\begin{quote}
\mr{काविळीत} \mr{डोळे} \mr{पिवळे} \mr{दिसतात}.
\end{quote}

\begin{quote}
\emph{kāviḷīt ḍoḷe pivḷe distāt.}
\end{quote}

\begin{quote}
In jaundice the eyes look yellow.
\end{quote}

Feminine.

Feminine; oblique \mr{काविळी}-: \mr{काविळीत}, in jaundice.

\begin{cousinweb}
\textbf{\mr{कावीळ}} (\emph{kāvīḷ}) is a tadbhava, inherited from Sanskrit \emph{kāmalā}; CDIAL lists Marathi \emph{kāvīḷ}, \textquotedblleft{}jaundice\textquotedblright{} (CDIAL 3044).
\end{cousinweb}

\begin{grammarlens}[title={What you've built}]
One more word for describing what you lived and felt.
\end{grammarlens}

\subsection*{Your turn}
\begin{itemize}
\raggedright
  \item \emph{Say it:} \emph{kāvīḷ}
  \item \emph{Say it:} \emph{kāvīḷ}, once more
  \item \emph{Say it:} \emph{kāviḷīt ḍoḷe pivḷe distāt.} — In jaundice the eyes look yellow.
  \item \emph{Recall:} say the Marathi for blood pressure, then the Marathi for asthma, then say \emph{kāvīḷ} again
\end{itemize}

\begin{tcolorbox}[breakable,colback=teal!4,colframe=teal!35!black,title={Before you move on}]
What does \mr{कावीळ} mean? (\textquotedblleft{}Jaundice\textquotedblright{}.) Say it once more.
\end{tcolorbox}
\section[\texorpdfstring{\mr{कांजिण्या} (kāñjiṇyā)}{कांजिण्या (kāñjiṇyā)}]{\mr{कांजिण्या} (kāñjiṇyā) — chickenpox}

\label{lesson:MR-C485-kanjinya}



\begin{quote}
Before the new one: say the Marathi for asthma, then the Marathi for jaundice.
\end{quote}

\subsection*{You'll want to know: \mr{कांजिण्या}}
\textbf{\mr{कांजिण्या}} — \emph{kāñjiṇyā} — \textquotedblleft{}chickenpox\textquotedblright{}.

\begin{quote}
\mr{लहानपणी} \mr{मला} \mr{कांजिण्या} \mr{झाल्या} \mr{होत्या}.
\end{quote}

\begin{quote}
\emph{lahānpaṇī malā kāñjiṇyā jhālyā hotyā.}
\end{quote}

\begin{quote}
I had chickenpox as a child.
\end{quote}

Feminine.

Plural and feminine: \mr{कांजिण्या} \mr{होणे}, with the dative.

\begin{cousinweb}
\textbf{\mr{कांजिण्या}} (\emph{kāñjiṇyā}) is a Marathi noun whose origin is not settled.
\end{cousinweb}

\begin{grammarlens}[title={What you've built}]
One more word for describing what you lived and felt.
\end{grammarlens}

\subsection*{Your turn}
\begin{itemize}
\raggedright
  \item \emph{Say it:} \emph{kāñjiṇyā}
  \item \emph{Say it:} \emph{kāñjiṇyā}, once more
  \item \emph{Say it:} \emph{lahānpaṇī malā kāñjiṇyā jhālyā hotyā.} — I had chickenpox as a child.
  \item \emph{Recall:} say the Marathi for asthma, then the Marathi for jaundice, then say \emph{kāñjiṇyā} again
\end{itemize}

\begin{tcolorbox}[breakable,colback=teal!4,colframe=teal!35!black,title={Before you move on}]
What does \mr{कांजिण्या} mean? (\textquotedblleft{}Chickenpox\textquotedblright{}.) Say it once more.
\end{tcolorbox}
